\documentclass{article}
\usepackage[utf8]{inputenc}
\usepackage[english]{babel}
\usepackage{amsmath}
\usepackage{amsthm}
\usepackage{amssymb}
\usepackage{color}
\usepackage{enumerate}
\usepackage{appendix}
\usepackage{algorithm}
\usepackage{algorithmic}
\usepackage{listings}
\usepackage{caption}
\usepackage{natbib}
\usepackage{graphicx}
\usepackage{multirow}
\usepackage{stmaryrd}


\newtheorem{theorem}{Theorem}
\newtheorem{lem}[theorem]{Lemma}
\newtheorem{cor}[theorem]{Corollary}
\newtheorem{rmk}{Remark}
\newtheorem{exe}{Example}
\newtheorem{prop}{Proposition}
\newtheorem{definition}{Definition}

\newcommand{\expc}[1]{\mathbb{E}\left[ #1 \right]}


\title{Notes on how to extend DeGAS}
\author{}
\date{}

\begin{document}

\section{Extension to extended Polynomials}

\subsection{Syntax Extension}

Let $x_1, \hdots, x_m$ be program variables.
We want to extend DeGAS' syntax to support assignments in the form:

$$ \mathtt{ x = \sum_{k=0}^K c_k m^{poly}_k \, \Bigg| \, \sum_{k=0}^K c_k m^{trig}_k \, \Bigg| \, \sum_{k=0}^K c_k m^{exp}_k}.$$

The previous include three different cases:
\begin{itemize}
    
    \item[i)] \textit{Polynomial assignments.} In this case $\mathtt{m^{poly}_k}$ are \textit{monomials}, defined as 
    $$ m^{poly} = x_1^{\alpha_1}...x_m^{\alpha_m}$$
    for $\alpha_i \in \mathbb{N}$.

    \item[ii)] \textit{Trigonometric polynomials.} In this case, $\mathtt{m^{trig}_k}$ are \textit{trigonometric monomials}, defined as:
    $$ m^{trig} = x_1^{\alpha_1}...x_{m}^{\alpha_m} \cos^{\beta_1}(x_1)...\cos^{\beta_m}(x_m)\sin^{\gamma_1}(x_1)...\sin^{\gamma_m}(x_m)$$
    for $\alpha_i, \beta_i, \gamma_i \in \mathbb{N}$.

    \item[iii)] \textit{Exponential polynomials.} In this case, $\mathtt{m^{exp}_k}$ are \textit{exponential monomials}, defined as:
    $$ m^{exp} = x_1^{\alpha_1}...x_m^{\alpha_m}\exp^{\beta_1}(x_1)...\exp^{\beta_m}(x_m)$$
    for $\alpha_i, \beta_i \in \mathbb{N}$.
   
\end{itemize}

Now we explain how DeGAS computes the semantics of these assignments.

\subsection{Semantics Extension}

\subsubsection{Reduction to monomials}

We show that to compute the means and the covariance matrix of $X_i$ after the assignment $\mathtt{x_i = \sum_{k=0}^K c_k m_k^*}$, where $* \in \{poly, trig, exp\}$ (but is always the same for all terms in the summation), it is sufficient to know how to compute the expectation of the single monomials.

For the mean we have
$$ \expc{X_i} = \expc{\sum_{k=0}^K c_k m_k^*} = \sum_{k=0}^K c_k \expc{m_k^*}  $$

For the covariance matrix e need to compute
\begin{align*} 
&Cov(X_i, X_j) = \expc{X_iX_j} - \expc{X_i}\expc{X_j} = \sum_{k=0}^K c_k \expc{x_jm_k^*} - \expc{X_i}\expc{X_j} \\
&Var(x_i) = \expc{X_i^2} - \expc{X_i}^2 = \sum_{k=0}^K c_k^2 \expc{(m_k^*)^2} + 2 \sum_{0 \le k_1 < k_2 \le K} c_{k_1} c_{k_2} \expc{m_{k_1}^*m_{k_2}^*}  - \expc{X_i}^2
\end{align*}
where $x_j m_k^*, (m_k^*)^2$ and $m_{k_1}^*m_{k_2}^*$ are all monomials of the same type as $m_k^*$.

\subsubsection{Expectation of monomials}

\paragraph{Monomials}

For computing $\expc{m_k^{poly}}$ we use the following.

\begin{theorem}[Isserlis' Theorem]
For a vector $X = (X_1, \hdots, X_n)$ which is jointly Gaussian, it holds:
$$ \expc{(X_1 - \mu_1)\hdots(X_n - \mu_n)} = \sum_{\substack{p = \text{partition of} \\ \{1, \hdots, m\} \text{into pairs} \\ \{(i_1, i_2), \hdots, (i_{n-1}, i_n)}\}} \prod_{(i_j, i_k) \in p} Cov(X_{i_j}, X_{i_k}) $$ 
\end{theorem}

From Isserlis' theorem one obtains a formulas for $\expc{X_1\hdots X_n}$ as a function of expectations of lower order monomials. In fact, at the l.h.s. one can compute the product and obtain a monic polynomial with highest order term $X_1\hdots X_n$. Then, applying the linearity of the expectation one directly obtains that $\expc{X_1\hdots X_n}$ can be expressed as a linear combination of the expectations of lower order monomials and the r.h.s. of the previous equation.

In general, for computing $\expc{X_1^{\alpha_1}\hdots X_m^{\alpha_k}}$ one has to set $Y_{1,i}=\hdots=Y_{1,\alpha_i}=X_i$ for $i=1, \hdots, m$. Then the vector $Y = (Y_{1,1}, \hdots, Y_{m,\alpha_m})$ is still jointly Gaussian and Isserlis' theorem can be applied to it.

\paragraph{Trigonometric Monomials}

For computing $\expc{m_k^{trig}}$ we use the following.

\begin{lem}
For a random vector $X = (X_1, \hdots, X_m)$ with characteristic function $\Phi_X(t)$:
\begin{align*}
    &\expc{X_1^{\alpha_1}\hdots X_m^{\alpha_m} \cos^{\beta_1}(X_1)\hdots\cos^{\beta_m}(X_m)\sin^{\gamma_1}(X_1)\hdots\sin^{\gamma_m}(X_m)} = \\
    &\frac{ \sum_{\substack{(b_1, \hdots, b_m)\\ = \\ (0, \hdots 0)}}^{\beta_1, \hdots, \beta_m} \sum_{\substack{(c_1, \hdots, c_m)\\ = \\  (0, \hdots, 0)}}^{\gamma_1, \hdots, \gamma_m} \prod_{i=1}^m \binom{\beta_i}{b_i} \binom{\gamma_i}{c_i} \frac{d^{\alpha_1}}{dt_1^{\alpha_1}} \hdots \frac{d^{\alpha_m}}{dt_m^{\alpha_m}} \Phi_X(t_1, \hdots, t_m)\Bigg|_{(t_i = 2(b_i + c_i) - \beta_i - \gamma_i)} }{2^{\sum_{i=1}^m (\beta_i + \gamma_i)}i^{\sum_{i=1}^m (\alpha_i + \gamma_i)}}
\end{align*}
\end{lem}

For a multivariate Gaussian $X = (X_1, \hdots, X_m)$ the characteristic function is (let $t = (t_1, \hdots, t_m)$ be a column vector):
$$\Phi_X(t) = \exp\left(i \mu^T \cdot t - \frac{1}{2} t^T \Sigma t\right)$$

\paragraph{Exponential Monomials}

For computing $\expc{m_k^{exp}}$ we use the following.

\begin{lem}
For a random vector $X = (X_1, \hdots, X_m)$ with moment-generating function $\mathcal{M}_X(t)$:
\begin{align*}
\expc{X_1^{\alpha_1}\hdots X_m^{\alpha_m}\exp^{\beta_1}(X_1)\hdots\exp^{\beta_m}(X_m)} = \frac{d^{\alpha_1}}{dt_1^{\alpha_1}} \hdots \frac{d^{\alpha_m}}{t_m^{\alpha_m}} \mathcal{M}_X(t_1, \hdots, t_m) \Bigg|_{(t_i = \beta_i)}
\end{align*}
\end{lem}

For a multivariate Gaussian $X = (X_1, \hdots, X_m)$ the moment-generating function is (let $t = (t_1, \hdots, t_m)$ be a column vector):
$$\mathcal{M}_X(t) = \exp\left(\mu^T \cdot t + \frac{1}{2} t^T \Sigma t\right)$$

\subsubsection{Generalization to Gaussian Mixtures}

In general DeGAS computes for Gaussian Mixtures, so the above results should be applied to every components of the mixture.

\subsection{Practical steps to implement this}

\begin{itemize}
    \item Modify the grammars in the files SOGA.g4 and ASGMT.g4 so that now assignment allows for polynomials, trigonometric polynomials and exponential polynomials
    \item Generate new files using antlr4
    \item Move the newly generated files to the src folder, and modify them with the additions specified in Snippet.txt 
    \item Modify the functions defined in libSOGAupdate.py so that now when a polynomial, a trigonometric polynomial or an exponential polynomial appears in an assignment the above formulas are used to compute the means and the covariance matrices of the transformed distributions.
\end{itemize}

Other practical notes:
\begin{itemize}
\item Use \texttt{x\textasciicircum n} for powers of x, \texttt{cos(x), sin(x), exp(x)} for sine, cosine and exponential of x respectively, and use \texttt{cos(x)\textasciicircum n, sin(x)\textasciicircum n} and \texttt{exp(x)\textasciicircum n} for powers of trig/exp.
\item Keep the current implementation for the linear case (probably more efficient). 
\item Note that in general when writing a monomial in a program the terms with exponent 0 will be omitted (so in a monomial not all variables will be explicitly written with their exponent). For terms with exponent 1 we do not write the exponent, but only the associated variable.
\item The trig and exp cases require computing derivative. Implement them analytically using Pytorch.
\item In general, \texttt{gm(...)} could appear inside a new monomial. 
\item We do NOT allow for mixed terms
\end{itemize}
\end{document}